% !TEX program = xelatex
% Lernplatt - by Can Siebert
% Kurs 71 (Dig 42) - Tag 5 - Aufgabenheft
% Von Diagrammen zu Enterprise-Faehigkeit: Visio/Lucidchart + EA + MDA

\documentclass[11pt,a4paper,oneside]{article}

\usepackage{polyglossia}
\setdefaultlanguage{german}
\usepackage[a4paper,margin=2.5cm]{geometry}

\usepackage{fontspec}
\defaultfontfeatures{Ligatures=TeX}
\IfFontExistsTF{Charter}{\setmainfont{Charter}}{%
  \setmainfont{Noto Serif}%
}
\IfFontExistsTF{Source Sans 3}{\setsansfont{Source Sans 3}}{%
  \IfFontExistsTF{Source Sans Pro}{\setsansfont{Source Sans Pro}}{\setsansfont{Liberation Sans}}%
}
\newcommand{\caveatfontfallback}{\itshape}
\IfFontExistsTF{Caveat}{\newfontfamily\caveatfont{Caveat}[Scale=1.15]}{\newcommand\caveatfont{\caveatfontfallback}}
\setmonofont{Liberation Mono}

\usepackage{microtype}
\usepackage{eurosym}
\linespread{1.15}

\usepackage{xcolor}
\definecolor{lpInk}{RGB}{20,28,38}
\definecolor{lpAccent}{RGB}{157,74,26}
\definecolor{lpAccentLight}{RGB}{246,228,212}
\definecolor{lpAmber}{RGB}{222,138,30}
\definecolor{lpAmberLight}{RGB}{255,243,217}
\definecolor{lpAmberBorder}{RGB}{196,118,18}
\definecolor{lpGray}{RGB}{96,104,114}
\definecolor{lpGrayLight}{RGB}{238,240,243}
\definecolor{lpGreen}{RGB}{34,120,72}
\definecolor{lpGreenLight}{RGB}{222,238,228}
\definecolor{lpGreenBorder}{RGB}{24,96,58}
\definecolor{lpL1}{RGB}{34,120,72}
\definecolor{lpL2}{RGB}{22,90,140}
\definecolor{lpL3}{RGB}{170,42,52}

\usepackage[most]{tcolorbox}
\tcbuselibrary{breakable,skins}

\newtcolorbox{infobox}[1][Hinweis]{%
  enhanced, breakable, colback=lpAccentLight, colframe=lpAccent,
  boxrule=0.6pt, arc=2pt, left=10pt,right=10pt,top=8pt,bottom=8pt,
  fonttitle=\sffamily\bfseries\color{white}, coltitle=white,
  title={#1}, attach boxed title to top left={xshift=8pt,yshift=-8pt},
  boxed title style={size=small, colback=lpAccent, colframe=lpAccent, boxrule=0.4pt, arc=1pt},
  top=14pt
}

\newtcolorbox{antwortbox}[1][Ihre Antwort]{%
  enhanced, breakable, colback=white, colframe=lpGray,
  boxrule=0.4pt, arc=2pt, left=10pt,right=10pt,top=8pt,bottom=8pt,
  fonttitle=\sffamily\bfseries\color{white}, coltitle=white,
  title={#1}, attach boxed title to top left={xshift=8pt,yshift=-8pt},
  boxed title style={size=small, colback=lpGray, colframe=lpGray, boxrule=0.4pt, arc=1pt},
  top=14pt
}

\usepackage{enumitem}
\setlist{itemsep=2pt, topsep=4pt, parsep=0pt}
\setlist[itemize,1]{label={\textbullet},leftmargin=1.6em}
\setlist[itemize,2]{label={\textendash},leftmargin=1.4em}
\setlist[enumerate,1]{label=\arabic*., leftmargin=1.8em}

\usepackage{tabularx}
\usepackage{array}
\usepackage{booktabs}
\usepackage{longtable}

\usepackage{fancyhdr}
\usepackage{lastpage}
\pagestyle{fancy}
\fancyhf{}
\renewcommand{\headrulewidth}{0.3pt}
\renewcommand{\footrulewidth}{0pt}
\fancyhead[L]{\footnotesize\sffamily\color{lpGray}Aufgabenheft \textperiodcentered{} Kurs 71 \textperiodcentered{} Tag 5}
\fancyhead[R]{\footnotesize\sffamily\color{lpGray}Visio/Lucidchart \textperiodcentered{} EA \textperiodcentered{} MDA}
\fancyfoot[L]{\footnotesize\sffamily\color{lpGray}\textcopyright{} Can Siebert}
\fancyfoot[R]{\footnotesize\sffamily\color{lpGray}\thepage{} / \pageref{LastPage}}

\fancypagestyle{plain}{%
  \fancyhf{}%
  \renewcommand{\headrulewidth}{0pt}%
  \fancyfoot[C]{\footnotesize\sffamily\color{lpGray}\thepage{} / \pageref{LastPage}}%
}

\usepackage[explicit]{titlesec}
\titleformat{\section}[block]
  {\sffamily\Large\bfseries\color{lpAccent}}
  {\thesection.}{0.6em}{#1}
  [{\vspace{2pt}\color{lpAccent}\titlerule[0.6pt]}]
\titleformat{\subsection}[block]
  {\sffamily\large\bfseries\color{lpInk}}
  {\thesubsection}{0.6em}{#1}
\titlespacing*{\section}{0pt}{14pt}{8pt}
\titlespacing*{\subsection}{0pt}{10pt}{4pt}

\usepackage[hidelinks,unicode]{hyperref}

\newcommand{\akzent}[1]{{\caveatfont\color{lpAmber}#1}}
\newcommand{\fachbegriff}[1]{\textbf{#1}}
\newcommand{\quelle}[1]{{\footnotesize\sffamily\color{lpGray}(#1)}}

\newcommand{\niveaubadge}[2]{%
  \tcbox[on line, colback=#1, colframe=#1, coltext=white,
         boxrule=0pt, arc=2pt, left=5pt,right=5pt,top=1pt,bottom=1pt,
         nobeforeafter]{\sffamily\bfseries\footnotesize #2}%
}
\newcommand{\nivLi}{\niveaubadge{lpL1}{L1\,\textbullet\,Wissen}}
\newcommand{\nivLii}{\niveaubadge{lpL2}{L2\,\textbullet\,Anwendung}}
\newcommand{\nivLiii}{\niveaubadge{lpL3}{L3\,\textbullet\,Management}}

\newcommand{\zeit}[1]{%
  \tcbox[on line, colback=lpGrayLight, colframe=lpGrayLight, coltext=lpInk,
         boxrule=0pt, arc=2pt, left=5pt,right=5pt,top=1pt,bottom=1pt,
         nobeforeafter]{\sffamily\footnotesize\textbf{$\sim$\,#1}}%
}

\newcommand{\aufgabenkopf}[4]{%
  \par\vspace{6pt}
  \noindent{\sffamily\Large\bfseries\color{lpAccent}Aufgabe #1 \textendash{} #2}\par
  \vspace{2pt}
  \noindent{\color{lpAccent}\rule{\linewidth}{0.6pt}}\par
  \vspace{4pt}
  \noindent #3 \hfill \zeit{#4}\par
  \vspace{6pt}
}

\newcommand{\ytquery}[1]{%
  \par\vspace{4pt}
  \noindent{\footnotesize\sffamily\color{lpGray}\textbf{Lernvideo zum Thema:}\ %
    \href{https://studyflix.de/search?query=#1}{\textcolor{lpAccent}{Studyflix-Treffer}}\ ·\ %
    \href{https://www.youtube.com/results?search\_query=#1}{\textcolor{lpAccent}{YouTube-Treffer}}\\%
    \textit{\textcolor{lpGray}{Stichworte: #1}}}\par
}

\newcommand{\antwortlinien}[1]{%
  \par\vspace{6pt}
  \foreach \x in {1,...,#1}{%
    \noindent\makebox[\linewidth]{\dotfill}\\[6pt]
  }
}
\usepackage{pgffor}

\begin{document}

\begin{titlepage}
  \pagestyle{empty}
  \vspace*{1.5cm}
  \noindent{\sffamily\footnotesize\color{lpGray}LERNPLATT \textperiodcentered{} BY CAN SIEBERT}\\[2pt]
  \noindent{\color{lpAccent}\rule{\linewidth}{1.2pt}}\\[4pt]
  \noindent{\sffamily\footnotesize\color{lpGray}AUFGABENHEFT \textperiodcentered{} MODUL 2 \textperiodcentered{} TAG 5 VON 16 \textperiodcentered{} KURS 71 (Dig 42) \textperiodcentered{} 3.\,Juni 2026}

  \vspace{3cm}
  {\sffamily\fontsize{30}{34}\selectfont\bfseries\color{lpInk}Aufgabenheft\\Tag 5\par}

  \vspace{0.7cm}
  {\sffamily\large\color{lpGray}Von Diagrammen zu Enterprise-\\F\"ahigkeit \textendash{} Visio/Lucidchart,\\EA-Rahmen und MDA.\par}

  \vspace{1.5cm}
  {\caveatfont\LARGE\color{lpAmber}11 Aufgaben \textperiodcentered{} L1 \textperiodcentered{} L2 \textperiodcentered{} L3}\par

  \vfill

  \noindent\begin{minipage}{0.55\linewidth}
    {\sffamily\footnotesize\color{lpGray}Niveau-Mix:}\\
    {\sffamily\small\color{lpInk}3\,\textsf{L1} \textperiodcentered{} 5\,\textsf{L2} \textperiodcentered{} 3\,\textsf{L3}}\\[10pt]
    {\sffamily\footnotesize\color{lpGray}Bearbeitung:}\\
    {\sffamily\small\color{lpInk}Selbstbearbeitung im Lernpfad}
  \end{minipage}\hfill
  \begin{minipage}{0.4\linewidth}
    \raggedleft
    {\sffamily\footnotesize\color{lpGray}Dozent:}\\
    {\sffamily\small\color{lpInk}Can Siebert}\\[10pt]
    {\sffamily\footnotesize\color{lpGray}Begleitend:}\\
    {\sffamily\small\color{lpInk}Lehrskript \& L\"osungsheft}
  \end{minipage}

  \vspace{0.5cm}
  \noindent{\color{lpAccent}\rule{\linewidth}{0.6pt}}
\end{titlepage}

\section*{So arbeiten Sie mit diesem Heft}
\thispagestyle{plain}

\noindent Dieses Aufgabenheft enth\"alt elf Aufgaben zu Tag 5 \textendash{} \emph{Von Diagrammen zu Enterprise-F\"ahigkeit}: kollaboratives Modellieren (Visio / Lucidchart), Enterprise Architecture (Capability Map) und Model-Driven Architecture (CIM / PIM / PSM). Die Aufgaben gliedern sich in drei Niveaus:

\begin{itemize}
  \item \nivLi{} \textendash{} \fachbegriff{Wissen.} Standards, Capability-Definition, MDA-Ebenen. 5\,--\,10 Minuten pro Aufgabe.
  \item \nivLii{} \textendash{} \fachbegriff{Anwendung.} Modelling Charter, Capability Map, MDA-Pfad, Naming-Audit, KPI-Set. 15\,--\,30 Minuten.
  \item \nivLiii{} \textendash{} \fachbegriff{Management.} Fallstudie \emph{RetailGroup Europe}, Tool-Strategie, Operating Model. 30\,--\,45 Minuten.
\end{itemize}

\noindent Jede Aufgabe enth\"alt:

\begin{itemize}
  \item eine Niveau-Markierung und einen Zeit-Richtwert,
  \item den Aufgabentext (oft mit Kontext oder Fallbeschreibung),
  \item einen Antwort-Freiraum f\"ur Ihre L\"osung,
  \item eine \emph{YouTube-Suchquery} f\"ur den begleitenden Lernpfad.
\end{itemize}

\begin{infobox}[Hinweis]
Die Musterl\"osungen finden Sie im \emph{L\"osungsheft Tag 5}. Versuchen Sie jede Aufgabe zuerst eigenst\"andig. Insbesondere auf L3 ist die \emph{Entscheidung} entscheidend, nicht die Vollst\"andigkeit der Beschreibung. Wer auf CIM-Ebene Plattform-Namen schreibt, hat MDA missverstanden \textendash{} und wer auf PSM-Ebene Business-Ziele formuliert, ebenfalls.
\end{infobox}

\clearpage

\section*{Niveau L1 \textendash{} Wissen \& Reproduktion}
\addcontentsline{toc}{section}{L1 -- Wissen}

% LERNPLATT_HILFSVIDEOS
\section*{Hilfsvideos zu diesem Tag}
\addcontentsline{toc}{section}{Hilfsvideos zu diesem Tag}
\thispagestyle{plain}

\noindent Die folgenden YouTube-Erkl\"arvideos vermitteln die Inhalte, die Sie zur Bearbeitung der Aufgaben ben\"otigen. Klicken Sie auf den Titel, um das Video direkt zu \"offnen; die vollst\"andige URL steht in Klammern.

\begin{description}[style=nextline,leftmargin=18pt,labelindent=0pt,itemsep=8pt]
  \item[1.~TOGAF in 12 Minuten — Enterprise Architecture in der Übersicht] \href{https://youtu.be/kp0yhrTKgg4}{\textcolor{lpAccent}{https://youtu.be/kp0yhrTKgg4}}\\[2pt]
  {\footnotesize\color{lpGray}(URL: \nolinkurl{https://youtu.be/kp0yhrTKgg4})}
  \item[2.~ArchiMate — Notation für Enterprise Architecture] \href{https://youtu.be/XGqbdfyuoaQ}{\textcolor{lpAccent}{https://youtu.be/XGqbdfyuoaQ}}\\[2pt]
  {\footnotesize\color{lpGray}(URL: \nolinkurl{https://youtu.be/XGqbdfyuoaQ})}
  \item[3.~Business Capability Mapping — was sind Capabilities?] \href{https://youtu.be/_cC76sI8QrM}{\textcolor{lpAccent}{\nolinkurl{https://youtu.be/_cC76sI8QrM}}}\\[2pt]
  {\footnotesize\color{lpGray}(URL: \nolinkurl{https://youtu.be/_cC76sI8QrM})}
  \item[4.~Zachman Framework in 9 Minuten — EA-Matrix demystifiziert] \href{https://youtu.be/zrA9Tia8k1Y}{\textcolor{lpAccent}{https://youtu.be/zrA9Tia8k1Y}}\\[2pt]
  {\footnotesize\color{lpGray}(URL: \nolinkurl{https://youtu.be/zrA9Tia8k1Y})}
\end{description}
\vspace{0.6em}
\noindent{\footnotesize\color{lpGray}\textit{Tipp:} \"Offnen Sie das Video, lassen Sie es im Hintergrund laufen und arbeiten Sie parallel an der Aufgabe.}

\clearpage

\aufgabenkopf{1}{Visio/Lucidchart als Enterprise-Werkzeug}{\nivLi}{8\,min}

\noindent Erkl\"aren Sie pr\"agnant: Was unterscheidet ein \emph{Enterprise-Modellierungs\-werkzeug} (z.\,B.\ Visio/Lucidchart in der Enterprise-Lizenz) von einer \emph{Einzeldatei} (z.\,B.\ eine \texttt{.vsdx} auf dem Laufwerk)? Listen Sie die \emph{vier Enterprise-Capabilities}:

\medskip
\noindent\textbf{Hinweis:} (a) \emph{Kollaboration} (gleichzeitiges Arbeiten, Kommentare), (b) \emph{Standards} (Templates, Shapesets, Naming), (c) \emph{Freigaben} (Lifecycle, Versionierung), (d) \emph{Single Source of Truth} (Verlinkung zu Policies, KPIs, Artefakten).

\medskip
\noindent F\"ur jede Capability:
\begin{enumerate}
  \item Beschreiben Sie in einem Satz, was sie konkret leistet.
  \item Geben Sie ein \emph{Anti-Muster} an, das ohne sie entsteht (``Ohne X passiert\ldots'').
\end{enumerate}

\ytquery{Lucidchart Visio Enterprise Kollaboration Shape Library Vergleich}

\noindent Erg\"anzen Sie eine Vergleichstabelle:

\medskip
\begin{tabularx}{\linewidth}{@{}p{4.4cm}|X|X@{}}
\toprule
\textbf{Capability} & \textbf{Mit Enterprise-Tool} & \textbf{Ohne (Einzeldatei)} \\
\midrule
Kollaboration       & & \\[12pt]
\midrule
Standards           & & \\[12pt]
\midrule
Freigaben           & & \\[12pt]
\midrule
Single Source of Truth & & \\[12pt]
\bottomrule
\end{tabularx}

\antwortlinien{8}

\clearpage

\aufgabenkopf{2}{Enterprise Architecture \textendash{} f\"unf Sichten}{\nivLi}{10\,min}

\noindent Enterprise Architecture (EA) beantwortet drei strategische Fragen:
\begin{itemize}
  \item \emph{Was k\"onnen wir?} \textendash{} Capabilities.
  \item \emph{Wie arbeiten wir?} \textendash{} Prozesse.
  \item \emph{Womit arbeiten wir?} \textendash{} Applikationen / Daten / Technologie.
\end{itemize}

\noindent Listen Sie die \emph{f\"unf typischen EA-Sichten} auf und beschreiben Sie pro Sicht: \emph{Was zeigt sie?}, \emph{Wer ist Zielgruppe?}, \emph{Welches Steuerungsproblem l\"ost sie?}

\medskip
\noindent\textbf{Hinweis:} Capability Map \textperiodcentered{} Prozesslandkarte \textperiodcentered{} Applikationslandschaft \textperiodcentered{} Datenobjekt-\"Ubersicht \textperiodcentered{} Schnittstellen-/Integrations-Map.

\ytquery{Enterprise Architecture Capability Map Sichten TOGAF ArchiMate Beispiel}

\antwortlinien{14}

\clearpage

\aufgabenkopf{3}{MDA \textendash{} CIM / PIM / PSM}{\nivLi}{8\,min}

\noindent Erkl\"aren Sie die drei Ebenen von \fachbegriff{Model-Driven Architecture (MDA)}:

\begin{enumerate}
  \item \emph{CIM (Computation Independent Model)} \textendash{} Business-Sicht: Ziele, Prozesse, Begriffe.
  \item \emph{PIM (Platform Independent Model)} \textendash{} logische System-/Service-Sicht (\emph{ohne} konkrete Plattform).
  \item \emph{PSM (Platform Specific Model)} \textendash{} konkrete technische Umsetzung (Plattform, Schnittstellen, Deployments).
\end{enumerate}

\medskip
\noindent F\"ur jede Ebene:
\begin{itemize}
  \item Wer ist die Zielgruppe / der Owner?
  \item Was \emph{geh\"ort hinein}, was \emph{nicht}?
\end{itemize}

\noindent Schlie{\ss}en Sie mit einem Satz: Warum ist die saubere Trennung CIM/PIM/PSM eine \emph{Verantwortungs}-Trennung \textendash{} und kein technisches Detail?

\ytquery{Model Driven Architecture MDA CIM PIM PSM Beispiel OMG}

\antwortlinien{10}

\clearpage

\section*{Niveau L2 \textendash{} Anwendung \& Transfer}
\addcontentsline{toc}{section}{L2 -- Anwendung}

\aufgabenkopf{4}{Modelling Charter in 45 Minuten}{\nivLii}{25\,min}

\begin{infobox}[Ausgangslage]
Sechs Teams modellieren Prozesse in unterschiedlichen Tools und Stilen. Management will Vergleichbarkeit, Audit will Versionen, Teams wollen Flexibilit\"at.
\end{infobox}

\noindent\textbf{Arbeitsauftrag:}

\begin{enumerate}
  \item \textbf{F\"unf nicht verhandelbare Standards:} z.\,B.\ Naming, Start/Ende, Rollen sichtbar, Status/Version, Owner. F\"ur jeden Standard: \emph{Begr\"undung} und \emph{Sanktion bei Versto{\ss}}.
  \item \textbf{Drei Freiheitsgrade:} z.\,B.\ Detailtiefe, Notation f\"ur \"Ubersichten, lokale Varianten. F\"ur jeden: warum bewusst freigegeben?
  \item \textbf{Freigabe- \& \"Anderungsprozess:} Draft \textrightarrow{} Review \textrightarrow{} Approved \textrightarrow{} Retired, inkl.\ Rollen (Modeler / Reviewer / Owner).
  \item \textbf{Entscheidung unter Unsicherheit:} Was setzen Sie \emph{sofort} um (2 Wochen), was erst \emph{sp\"ater} (3 Monate) \textendash{} mit Begr\"undung.
\end{enumerate}

\noindent Abgabe: 1-Seiten-\emph{Modelling Charter}.

\noindent\textbf{Anti-Beispiel zum Reflektieren:} Eine Charter, die 24 Regeln enth\"alt, wird zu Tapete. Eine, die 5 Regeln mit Z\"ahnen hat, steuert. Was machen Sie, wenn ein Fachbereich-Lead sagt: ``Bei uns funktioniert das anders''? Notieren Sie eine Linie.

\ytquery{Modelling Charter Process Standards Naming Convention BPM Governance}

\antwortlinien{20}

\clearpage

\aufgabenkopf{5}{Capability Map ``Digital Customer Onboarding''}{\nivLii}{22\,min}

\begin{infobox}[Mini-Fall]
Ein Unternehmen will \emph{Digital Customer Onboarding} verbessern: langsam, viele R\"uckfragen, Compliance-Risiko, IT-Schnittstellen br\"uchig.
\end{infobox}

\noindent\textbf{Arbeitsauftrag:}

\begin{enumerate}
  \item \textbf{Capability Map} mit \emph{8\,--\,12 Capabilities} (z.\,B.\ ``Kunde identifizieren'', ``Vertrag pr\"ufen'', ``Zugang provisionieren''). Optional in 2\,--\,3 Cluster.
  \item Markieren Sie die \emph{Top-3 Capabilities mit h\"ochstem Hebel} auf Durchlaufzeit oder Compliance.
  \item Zeichnen Sie eine einfache \emph{Prozesslandkarte} (E2E) und ordnen Sie Capabilities zu (welche Capability wird wo gebraucht?).
  \item Definieren Sie \emph{3 kritische Datenobjekte} (z.\,B.\ Kunde, Vertrag, Identit\"atsnachweis) und \emph{2 Risiken} (Datenqualit\"at, Datenschutz).
\end{enumerate}

\noindent\textbf{Bonus-Frage:} Welche \emph{eine Capability} bleibt in Ihrem Cluster bewusst ohne Hebel-Markierung \textendash{} und warum? (Senior-Reflex: nicht alles ist gleich wichtig.)

\ytquery{Capability Map Customer Onboarding Enterprise Architecture Hebel}

\antwortlinien{20}

\clearpage

\aufgabenkopf{6}{MDA-Pfad f\"ur eine priorisierte Capability}{\nivLii}{25\,min}

\begin{infobox}[Vorgabe]
W\"ahlen Sie eine der priorisierten Capabilities aus Aufgabe 5 (z.\,B.\ \emph{``Identit\"atspr\"ufung''} oder \emph{``Zugang provisionieren''}). Sie f\"uhren den MDA-Pfad CIM \textrightarrow{} PIM \textrightarrow{} PSM durch.
\end{infobox}

\noindent\textbf{Arbeitsauftrag:}

\begin{enumerate}
  \item \textbf{CIM:} Formulieren Sie das \emph{Business-Ziel} (z.\,B.\ ``Kunde in $<$\,24\,h aktiv''), den \emph{Prozess-Scope} und die Definition des zentralen Datenobjekts \emph{Kunde} (Pflichtfelder).
  \item \textbf{PIM:} Identifizieren Sie \emph{3 logische Services} (Name, Input, Output) \emph{ohne} konkrete Plattform. Skizzieren Sie den Datenfluss zwischen Services und 2 Fehlerf\"alle.
  \item \textbf{PSM:} Skizzieren Sie die konkrete Umsetzung: Welche bestehende Plattform (IAM, Workflow-Engine), welche Schnittstellen, welche Audit-Logging-Pflicht?
  \item Schlie{\ss}en Sie mit \emph{einem Satz}: Welche Entscheidung gilt auf CIM-Ebene \emph{nicht verhandelbar} \textendash{} und welche darf auf PSM-Ebene noch ge\"andert werden?
\end{enumerate}

\ytquery{MDA CIM PIM PSM Beispiel Identity Service Onboarding Capability}

\antwortlinien{18}

\clearpage

\aufgabenkopf{7}{Naming \& Konvention \textendash{} 8 Mini-F\"alle}{\nivLii}{15\,min}

\begin{infobox}[Acht reale Modellname-Beispiele]
\textbf{A.} \texttt{O2C\_final\_v3\_neu.vsdx}\\
\textbf{B.} \texttt{Bestellung}\\
\textbf{C.} \texttt{Order-to-Cash}\\
\textbf{D.} \texttt{ProcRevAB\_Hoehl\_2025.vsdx}\\
\textbf{E.} \texttt{Rechnung erstellen v1.0 - Draft - DE}\\
\textbf{F.} \texttt{Onboarding (cleaner version - DO NOT DELETE)}\\
\textbf{G.} \texttt{P2P Subprozess - QA - 2026-06-02}\\
\textbf{H.} \texttt{Untitled Diagram (37)}
\end{infobox}

\noindent\textbf{Arbeitsauftrag:}

\noindent F\"ur jedes der acht Beispiele:
\begin{enumerate}
  \item Bewerten Sie auf einer Skala \emph{gut / mittel / schlecht}.
  \item Begr\"unden Sie in einem Satz.
  \item Schlagen Sie \emph{einen korrigierten Namen} vor (nach Modelling-Charter-Standard: Prozess-/Subprozess, Status, Version, Standort, ohne pers\"onliche Marker).
\end{enumerate}

\noindent\textbf{Hinweis:} Es geht nicht darum, ``sch\"on'' zu benennen, sondern \emph{rezeptierbar}: Wer das Modell sucht, soll \emph{ohne Insider-Wissen} finden, ob es aktuell ist, wem es geh\"ort und f\"ur welchen Standort es gilt.

\ytquery{Naming Convention BPM Files Best Practice Versioning Process Model}

\antwortlinien{18}

\clearpage

\aufgabenkopf{8}{KPI-Set f\"ur Modellqualit\"at}{\nivLii}{18\,min}

\begin{infobox}[Vorgabe]
Eine Modelling Charter ohne Steuerung verkommt zur Papiertiger-Sammlung. Definieren Sie ein KPI-Set, das \emph{Modellqualit\"at} und \emph{Nutzen} messbar macht.
\end{infobox}

\noindent\textbf{Arbeitsauftrag:}

\begin{enumerate}
  \item \textbf{Vier KPIs}: zwei \emph{Outcome / Speed} (Wirkung des Repositories), zwei \emph{Risk / Compliance}. Pro KPI: Name, Formel, Quelle, Frequenz, Owner.
  \item \textbf{Pro KPI ein Manipulationsrisiko:} Wie k\"onnte sie ``gespielt'' werden?
  \item \textbf{Kopplung (Goodhart-Schutz):} Wie verhindern Sie, dass Speed-KPIs Qualit\"at zerst\"oren?
  \item \textbf{Entscheidung:} Welche \emph{eine KPI} f\"uhren Sie nicht ein \textendash{} und warum (Akzeptanzproblem, Datenrisiko)?
\end{enumerate}

\noindent Erg\"anzen Sie eine \emph{Schluss\-bemerkung} (3 S\"atze): Wie messen Sie, ob Ihr KPI-Set \emph{Wirkung} hat \textendash{} und nicht nur \emph{Reporting} erzeugt?

\ytquery{KPI Modellqualitaet Process Repository Goodhart Compliance Steuerung}

\antwortlinien{18}

\clearpage

\section*{Niveau L3 \textendash{} Management \& Entscheidungsarchitektur}
\addcontentsline{toc}{section}{L3 -- Management}

\aufgabenkopf{9}{Fallstudie \emph{RetailGroup Europe} \textendash{} Enterprise Model Governance + EA/MDA Roadmap}{\nivLiii}{45\,min}

\begin{infobox}[Fallstudie \emph{RetailGroup Europe}]
\textbf{Unternehmen:} \emph{RetailGroup Europe} (Filialen + Online).\\
\textbf{Problem:} Viele Prozessdokumente, keine Vergleichbarkeit; Digitalisierung stockt, weil Business und IT aneinander vorbeireden.

\smallskip
\textbf{Restriktionen:} \textbf{10 Wochen} bis Audit-Vorbereitung, Budget \textbf{100\,k\,\euro}, zwei parallel laufende IT-Projekte, unklare Verantwortlichkeiten.
\end{infobox}

\noindent\textbf{Arbeitsauftrag:}

\begin{enumerate}
  \item \textbf{Modelling Charter}: Standards + Freiheiten + Freigabeprozess.
  \item \textbf{Capability Map} (10\,--\,12) und Wahl von \emph{2 Capabilities} als Digitalisierungs\-priorit\"at.
  \item Definieren Sie f\"ur \emph{eine} priorisierte Capability den \textbf{MDA-Pfad}:
    \begin{itemize}
      \item CIM: Business-Ziel + Prozess-Scope + Datenobjekt-Pflichtfelder.
      \item PIM: Service-/Datenlogik, Fehlerf\"alle.
      \item PSM: konkrete Umsetzungsidee.
    \end{itemize}
  \item \textbf{4 KPIs} (2 Outcome / Speed, 2 Risk / Compliance) inkl.\ Owner.
  \item \textbf{Management\-entscheidung:} Welche \emph{2 Quick Wins} in den ersten 4 Wochen? Welche 2 strukturellen Ma{\ss}nahmen f\"ur Wochen 5\,--\,10?
\end{enumerate}

\noindent\textbf{Format-Hinweis:} Pr\"asentieren Sie das Ergebnis als \emph{eine A4-Seite}: oben Modelling Charter (Top-5 Bullets), Mitte Capability Map (Skizze), unten MDA-Pfad (Tabelle) und KPI-Set. ``Senior-Format'': es muss in 90 Sekunden lesbar sein.

\ytquery{Enterprise Model Governance Capability Map MDA Roadmap Retail Audit}

\antwortlinien{36}

\clearpage

\aufgabenkopf{10}{Tool-Strategie: Standard vs.\ f\"oderiert}{\nivLiii}{25\,min}

\noindent\textbf{Diskussionsaufgabe.}

\noindent \emph{RetailGroup} steht vor einer Tool-Entscheidung: \emph{ein Standard\-tool} f\"ur alle Bereiche, oder \emph{f\"oderierte} Toollandschaft mit klaren \"Ubergangs-Regeln. Beide Optionen haben gravierende Folgen.

\medskip
\noindent\textbf{Arbeitsauftrag:}

\begin{enumerate}
  \item \textbf{Trade-off-Matrix}: Stellen Sie beide Optionen in 4 Dimensionen gegen\"uber: \emph{Kollaboration}, \emph{Vergleichbarkeit}, \emph{Akzeptanz}, \emph{Migrations\-aufwand}.
  \item \textbf{Entscheidung}: Welche Option w\"ahlen Sie? Begr\"undung in 4\,--\,5 S\"atzen.
  \item \textbf{Risiko und Fr\"uhwarnsignal}: Welches Risiko nehmen Sie bewusst in Kauf? Welches Signal w\"urde Sie zum Wechsel zwingen?
  \item \textbf{Migrationspfad} (bei f\"oderiert \textrightarrow{} Standard): Welche 3 Schritte sind die \emph{unumgehbaren} Voraussetzungen?
\end{enumerate}

\ytquery{BPM Tool Strategy Standard vs Federated Visio Lucidchart Enterprise Decision}

\antwortlinien{18}

\clearpage

\aufgabenkopf{11}{Transferprojekt \textendash{} Enterprise Modelling Operating Model}{\nivLiii}{45\,min}

\begin{infobox}[Rolle und Ausgangslage]
\textbf{Ihre Rolle:} Chief Enterprise Architect / Chief Transformation Officer / Head of Governance.

\smallskip
\textbf{Auftrag:} In \textbf{12 Wochen} ein unternehmensweites Modellierungs- und Architektur-Betriebs\-modell etablieren, das Projekte beschleunigt, Auditf\"ahigkeit erh\"oht und Verantwortlichkeiten kl\"art.

\smallskip
\textbf{Restriktionen:} unterschiedliche Tool-Reife, Widerstand aus Fachbereichen, Budget \textbf{180\,k\,\euro}, parallele IT-Programme, unklare Datenverantwortung.
\end{infobox}

\noindent\textbf{Arbeitsauftrag:}

\begin{enumerate}
  \item \textbf{Entscheidungs\-architektur:} \emph{Top-8 Entscheidungen}, die durch Modelle / EA unterst\"utzt werden (Priorisierung, Standardisierung, Investitionen, Risiken, Schnittstellen).
  \item \textbf{Modelling Charter v2:} nicht verhandelbare Standards, Freiheitsgrade, Freigabe-/Change-Prozess, KPI-Set f\"ur Modell\-qualit\"at und Nutzen.
  \item \textbf{EA-Minimum Viable Architecture}: Capability Map + Prozesslandkarte + Applikations-/Datenobjekt-Minimum inkl.\ Verantwortliche (Data Owners).
  \item \textbf{MDA-Guardrails:} Was muss CIM enthalten (Definitionen, Ziele, Scope), was geh\"ort ins PIM, wann PSM \textendash{} inkl.\ Review-Gates.
  \item \textbf{Rollout \& Enablement:} 3 Wellen (Pilot \textrightarrow{} Scale \textrightarrow{} Sustain), Trainingspfad Level 1\,--\,3, Akzeptanzstrategie.
  \item \textbf{Zwei Entscheidungen unter Unsicherheit:}
    \begin{itemize}
      \item \emph{Tool-Strategie:} ``Standardtool'' vs.\ ``f\"oderiert'' \textendash{} Entscheidung + Risiko + Fr\"uhwarnsignal.
      \item \emph{Governance-Strenge:} minimal vs.\ streng \textendash{} Entscheidung + Akzeptanzplan + Trigger zum Nachsch\"arfen.
    \end{itemize}
\end{enumerate}

\noindent\textbf{Bewertungskriterien:} Strategie-Anschluss \textperiodcentered{} Verantwortungs\-klarheit \textperiodcentered{} Pragmatismus unter Restriktionen \textperiodcentered{} Entscheidungslogik und Risikoreife.

\noindent\textbf{Senior-Reflex:} Eine reine Ma{\ss}nahmenliste ist kein Operating Model. Sie m\"ussen erkl\"aren, \emph{warum} diese Entscheidungen in dieser Reihenfolge fallen, \emph{wer} sie tr\"agt und \emph{woran} Sie erkennen, dass es nicht funktioniert.

\ytquery{Enterprise Modelling Operating Model EA MDA Governance Rollout CEA}

\antwortlinien{38}

\clearpage

\section*{Abschluss}
\thispagestyle{plain}

\noindent Sie haben alle elf Aufgaben bearbeitet. Vergleichen Sie nun Ihre Antworten mit dem \emph{L\"osungsheft Tag 5}. Pr\"ufen Sie insbesondere folgende drei Punkte:

\begin{itemize}
  \item Habe ich die Modelling Charter \emph{tats\"achlich} verwendet, um eine Tool-Entscheidung zu treffen?
  \item Sind meine Capabilities wirklich F\"ahigkeiten \textendash{} oder Prozesse mit anderem Namen?
  \item Habe ich CIM, PIM und PSM sauber getrennt \textendash{} oder Business-Anforderungen mit Technik-Details verschmolzen?
  \item Habe ich beim Naming-Audit die acht Beispiele \emph{nach Konvention} bewertet \textendash{} oder nach pers\"onlichem Geschmack?
  \item Hat mein KPI-Set einen \emph{wirksamen} Goodhart-Schutz \textendash{} oder ist die Kopplung kosmetisch?
\end{itemize}

\medskip
\begin{infobox}[Was Sie jetzt k\"onnen sollten]
\begin{itemize}
  \item Enterprise-Modellierungs\-werkzeuge (Visio / Lucidchart) als \emph{Kollaborationsplattform} statt Zeichenprogramm einsetzen.
  \item Capability Map und Prozesslandkarte als zwei unabh\"angige EA-Sichten verstehen und kombinieren.
  \item Den MDA-Pfad CIM \textrightarrow{} PIM \textrightarrow{} PSM als \emph{Verantwortungs}-Trennung nutzen.
  \item Eine Modelling Charter so schreiben, dass sie Akzeptanz erzeugt \textendash{} statt Widerstand.
  \item Eine Tool-Strategie (Standard vs.\ f\"oderiert) auf Basis von Trade-offs entscheiden.
  \item Naming und Lifecycle-Status in der Datei-Benennung sichtbar machen \textendash{} damit Audit, Owner und Aktualit\"at recherchierbar sind.
  \item Modellqualit\"at messbar machen \textendash{} mit Goodhart-Schutz \"uber gekoppelte KPIs.
\end{itemize}
\end{infobox}

\medskip
\noindent\textbf{Hinweis f\"ur Tag 6:} Morgen verlassen wir die Modellierungs-Ebene und gehen in die \emph{Ausf\"uhrung} \textendash{} BPMS-Plattformen, API-Vertr\"age, Cloud-native Skalierung. Bringen Sie Ihre Capability-Priorit\"aten von heute mit \textendash{} sie sind die Grundlage f\"ur die Plattform-Entscheidung von morgen.

\vspace{1cm}
\noindent{\color{lpAccent}\rule{\linewidth}{0.6pt}}\\[4pt]
{\footnotesize\sffamily\color{lpGray}Lernplatt \textperiodcentered{} by Can Siebert \textperiodcentered{} Kurs 71 (Dig 42) \textperiodcentered{} Tag 5 \textperiodcentered{} Aufgabenheft \textperiodcentered{} \textcopyright{} 2026}

\end{document}
